\begin{abstract}
A fundamental tension exists in our understanding of multi-agent large language model (LLM) systems: in controlled laboratory experiments, LLM populations spontaneously develop shared behavioral conventions—turn-taking norms, information-sharing protocols, collective decision-making patterns—yet in open, deployed environments, no such emergence has been reliably observed \citep{ashery2025,li2026tsinghua}. This contradiction—the \textbf{Moltbook Illusion}—raises a critical question: \emph{are social norms an intrinsic property of LLM populations, or an artifact of laboratory environmental design?}

We present a systematic answer through a fully crossed $2\times2\times2\times2$ factorial experiment ($N=480$ runs) varying four environmental factors independently and in combination. We find that \emph{communication structure} is the dominant determinant of norm emergence (largest effect, $\eta^2=0.31$, $p<0.001$), while structured feedback loops are necessary but not sufficient for norm crystallization. Critically, our ``wild-type'' conditions reproduce the 93\% non-response inertia rate observed in Tsinghua Moltbook deployment data \citep{li2026tsinghua}, validating ecological fidelity.

We introduce the \textbf{Norm Emergence Environmental Model (NEEM)}: a framework predicting norm emergence as a function of environmental configuration rather than model capability. NEEM's central thesis is that norm emergence is an infrastructure outcome—where you place the feedback, who can talk to whom, and whether the population stays stable—determines whether a population develops shared behavioral expectations. Our findings show that communication infrastructure—guaranteed message delivery, bounded latency, shared context windows—produces norm emergence at low cost, while model upgrades alone do not.

\end{abstract}

\vspace{4pt}
\textbf{Keywords:} norm emergence, multi-agent LLMs, social conventions, environmental design, HCI, AI ethics
